\section*{数据与代码可用性}
\addcontentsline{toc}{section}{数据与代码可用性}

案例数据来自公开的 DeepCRISPR 数据集\citep{Chuai2018GenomeBiol}；处理后特征位于 \texttt{data/processed/}，数据划分与泄漏审计代码位于 \texttt{core/data/splitting/}。

本文全部定量结论仅来自经过验收的权威批次 \texttt{ultimate\_run}（\texttt{results/batches/ultimate\_run/}）：1\,344 次运行，`single`/`all`/`mixed` 各 448。该批次的逐 run 验收明细（覆盖度、划分摘要复核、序列与反向互补重叠、数据/代码/环境指纹、设备）由 \texttt{deploy/hpc/verify\_hpc\_rerun.py} 生成并保存为 \texttt{results/tables/audit/ultimate\_run\_verify\_\{runs,summary,report\}.csv/md}。存在已确认划分泄漏的历史批次 \texttt{batch\_20260909\_full} 已标记废弃，仅用于 \S\ref{sec:res-leakage} 的前后对照，不作为任何结论依据。

图表与表格由结果文件重新计算生成，不手工填写：\texttt{analysis/reporting/paper/make\_assets.py --batch ultimate\_run}（图 1--7 与表 1--6）、\texttt{analysis/reporting/paper\_analysis/build\_paper\_tables.py --batch ultimate\_run}（论文中间表）、\texttt{analysis/reporting/paper/make\_data\_quality\_table.py --batch ultimate\_run}（表~\ref{tab:dataquality}）。分析方案（AnalysisPlan）、任务执行状态与每个任务的产物路径随批次一并保存（\texttt{summary/analysis\_plan.json}、\texttt{analysis\_status.json}、\texttt{execution\_log.json}），可逐项追溯（附录~\ref{app:provenance}）。

\section*{作者贡献}
\addcontentsline{toc}{section}{作者贡献}
（待补）

\section*{利益冲突声明}
\addcontentsline{toc}{section}{利益冲突声明}
作者声明不存在利益冲突。
